\documentclass[12pt]{article}
\usepackage{amsmath,amssymb,amsthm}
\usepackage{geometry}\geometry{margin=1in}
\usepackage{hyperref}
\usepackage{booktabs}

\newtheorem{conjecture}{Conjecture}
\newtheorem{proposition}[conjecture]{Proposition}
\theoremstyle{definition}
\newtheorem{remark}[conjecture]{Remark}

\newcommand{\ph}{\varphi}
\newcommand{\logph}{\log\ph}
\newcommand{\Hstar}{\mathcal{H}_*}

\title{Golden Ratio Quantisation of Gravitational Wave Peak Frequencies\\
and Dark Matter Mass from the $\ph$-Lattice}
\author{Marvin L.\ Gentry\\
\small Independent Researcher, Seattle, WA\\
\small \texttt{drmlgentry@protonmail.com}\quad
ORCID: 0009-0006-4550-2663}
\date{April 2026}

\begin{document}
\maketitle

\begin{abstract}
The gravitational wave-induced freeze-in mechanism of Maleknejad and
Kopp~\cite{MaleknejadKopp2026} shows that a stochastic gravitational
wave (GW) background in the early Universe can produce massless Weyl
fermions which, if they later acquire mass, constitute dark matter.
We observe that this mechanism provides a natural cosmological
realisation of the golden ratio lattice $m = m_e\ph^{q/4}$ that
governs the Standard Model fermion mass spectrum in the Hyperbolic
Flavor Geometry (HFG) framework.
We conjecture that the GW peak frequency is quantised in units of
$\log\ph$: $q_{\rm peak}/\Hstar = \ph^k$ for integer $k$, and that
the produced dark matter mass satisfies $M_{\rm DM} = m_e\ph^{q/4}$.
The relic density constraint then relates $k$ and $q$, yielding
falsifiable predictions for future GW detectors (Einstein Telescope,
Cosmic Explorer) and direct dark matter searches.
\end{abstract}

\section{Introduction}

The Hyperbolic Flavor Geometry (HFG) framework~\cite{Gentry:Lucas,
Gentry:SU} establishes that the fermion mass spectrum obeys
\begin{equation}
  m = m_e \cdot \ph^{q/4},\qquad q \in \tfrac{1}{4}\mathbb{Z},
  \label{eq:lattice}
\end{equation}
where $\ph=(1+\sqrt5)/2$ and $m_e$ is the electron mass.
The Lucas sequence $L_k = \ph^k + \ph^{-k}$ governs the covering
tower primes of the flavor manifolds, and $\log\ph$ is the
fundamental geodesic length unit.

Maleknejad and Kopp~\cite{MaleknejadKopp2026} recently showed that
a stochastic GW background breaks the conformal symmetry of massless
Weyl fermions, producing them via a 1-loop graviton--fermion vertex.
The produced fermion energy density scales as~\cite{MaleknejadKopp2026}
\begin{equation}
  \Omega_{\rm DM,0} \approx 0.36\,\mathcal{C}\,
  \left(\frac{g_*}{106.75}\right)^{\!4/3}
  \!\!\left(\frac{\Omega_{\rm peak}}{10^{-6}}\right)
  \left(\frac{M}{T_*}\right)
  \left(\frac{q_{\rm peak}/\Hstar}{100}\right)^{\!4}
  \!\!\left(\frac{T_*}{3\times10^{11}\,{\rm GeV}}\right)^{\!5},
  \label{eq:relic}
\end{equation}
where $\mathcal{C}=\mathcal{O}(1)$ depends on the GW spectral index
and coherence, $T_*$ is the production temperature, $\Hstar$ is the
comoving Hubble scale at $T_*$, and $\Omega_{\rm peak}$ is the peak
GW spectral energy density.
Here we point out that equation~\eqref{eq:relic} can accommodate
the $\ph$-lattice in a natural and falsifiable way.

\section{$\ph$-lattice quantisation}

\subsection{Peak frequency quantisation}

The natural scale for any logarithmically structured spectrum in the
HFG framework is $\log\ph$, the regulator of $\mathbb{Q}(\sqrt5)$
and the fundamental geodesic unit of the flavor
manifolds~\cite{Gentry:Lucas}.
We conjecture:

\begin{conjecture}[$\ph$-quantised GW spectrum]
\label{conj:gw}
The peak frequency of a primordial stochastic GW background satisfies
\[
  \frac{q_{\rm peak}}{\Hstar} = \ph^k, \qquad k \in \mathbb{Z}_{\ge 0},
\]
i.e., $\log(q_{\rm peak}/\Hstar) = k\log\ph$.
\end{conjecture}

Typical GW sources give $q_{\rm peak}/\Hstar \sim 10$--$10^3$ for
first-order phase transitions~\cite{Caprini2016}, corresponding to
$k \sim 5$--$15$ (since $\ph^{10}\approx 123$, $\ph^{15}\approx 1364$).

\subsection{Dark matter mass on the $\ph$-lattice}

If the produced fermions acquire mass via a Higgs mechanism, the HFG
framework predicts $M_{\rm DM} = m_e\ph^{q/4}$ for some
$q\in\frac{1}{4}\mathbb{Z}$ not in the visible sector
Lucas set~\cite{Gentry:SU}.
Representative dark sector candidates from non-Lucas covering tower
primes are listed in Table~\ref{tab:dm}.

\begin{table}[htbp]
\centering
\caption{Dark sector mass candidates from non-Lucas covering primes.}
\label{tab:dm}
\begin{tabular}{ccll}
\toprule
Prime $p$ & $q=p$ & $M = m_e\ph^{p/4}$ & Note \\
\midrule
13 & 13 & 1.85 MeV  & sterile neutrino range \\
19 & 19 & 6.4 MeV   & light dark matter \\
31 & 31 & 39 MeV    & dark sector mediator \\
41 & 41 & 580 MeV   & sub-GeV dark matter \\
43 & 43 & 940 MeV   & near proton mass \\
\bottomrule
\end{tabular}
\end{table}

\subsection{Consistency relation and falsifiability}

Substituting Conjecture~\ref{conj:gw} and $M=m_e\ph^{q/4}$ into
equation~\eqref{eq:relic} with $\Omega_{\rm DM,0}=0.27$ gives a
constraint on the integers $k$ and $q$:
\begin{equation}
  \frac{m_e}{T_*}\,\ph^{4k + q/4} \approx
  \frac{0.75}{\mathcal{C}}
  \left(\frac{106.75}{g_*}\right)^{\!4/3}
  \!\!\left(\frac{10^{-6}}{\Omega_{\rm peak}}\right)
  \left(\frac{T_*}{3\times10^{11}\,{\rm GeV}}\right)^{\!-5}.
  \label{eq:constraint}
\end{equation}
For each choice of $(T_*, \Omega_{\rm peak}, \mathcal{C})$,
equation~\eqref{eq:constraint} selects a discrete set of
integer pairs $(k,q)$.

\begin{proposition}[Falsifiable prediction]
For a GW background detected with peak frequency $f_{\rm peak}$
(in Hz today), the ratio $\log(f_{\rm peak}/H_0)/\log\ph$ should be
close to an integer $k$. The dark matter particle produced by this
mechanism must have mass $M_{\rm DM}/m_e = \ph^{q/4}$ for some $q$
not in the Lucas set. Both conditions can be tested independently.
\end{proposition}

\begin{remark}[The 17 MeV anomaly]
The claimed 17 MeV dark photon anomaly implies $q/4 \approx 7.28$,
i.e., $q \approx 29.1 \approx L_7$ (a Lucas prime). The HFG
framework therefore predicts this anomaly is \emph{visible sector},
not dark sector — independently of the GW mechanism.
\end{remark}

\section{Discussion}

The stochastic GW background acts as a ``cosmic ruler'' imprinting
the $\ph$-lattice on the dark matter mass through the production
mechanism. The lattice would then appear in three observationally
distinct places:
\begin{enumerate}
  \item The GW peak frequency ratio $q_{\rm peak}/\Hstar = \ph^k$
        (detectable by Einstein Telescope~\cite{ET} or
        Cosmic Explorer~\cite{CE} in the kHz--GHz band).
  \item The dark matter mass $M_{\rm DM} = m_e\ph^{q/4}$
        (testable by direct detection or collider experiments).
  \item The visible sector fermion masses already confirmed by the
        HFG framework~\cite{Gentry:Lucas}.
\end{enumerate}

The Maleknejad--Kopp mechanism provides the first proposed \emph{physical
mechanism} by which the $\ph$-lattice could be imprinted cosmologically,
as opposed to being postulated as a symmetry of the low-energy effective
theory. If the GW background itself arises from a phase transition
governed by the same hyperbolic geometry (e.g., through the icosahedral
symmetry $A_5$ shared by the quasicrystal phonon spectrum and the HFG
flavor manifolds), the quantisation of $q_{\rm peak}$ would follow
naturally.

\section{Conclusion}

The gravitational wave-induced freeze-in mechanism provides a
cosmological arena in which the $\ph$-lattice of the HFG framework
can be tested. The conjecture $q_{\rm peak}/\Hstar = \ph^k$ and the
dark matter mass constraint $M_{\rm DM} = m_e\ph^{q/4}$ together
with the relic density equation~\eqref{eq:relic} yield a discrete,
falsifiable prediction. Confirmation by future GW detectors would
establish the $\ph$-lattice as a fundamental structure of nature
beyond the Standard Model flavor sector.

\section*{Data Availability}
No new data were generated.

\section*{Conflict of Interest}
The author declares no conflicts of interest.

\begin{thebibliography}{99}
\bibitem{MaleknejadKopp2026}
A.~Maleknejad and J.~Kopp,
\textit{Gravitational-Wave Induced Freeze-In of Fermionic Dark Matter},
Phys.\ Rev.\ Lett.\ \textbf{136}, 131501 (2026).
\href{https://doi.org/10.1103/lr69-45v8}{doi:10.1103/lr69-45v8}

\bibitem{Gentry:Lucas}
M.~L.~Gentry,
\textit{The Lucas Sequence as the Arithmetic Bridge in Hyperbolic Flavor Geometry},
J.\ Number Theory, submitted (2026). JNTH-D-26-00538.
\href{https://doi.org/10.2139/ssrn.6660598}{SSRN:6660598}

\bibitem{Gentry:SU}
M.~L.~Gentry,
\textit{Spectral Universality and Slope-Encoded Arithmetic in Dehn Fillings of Hyperbolic 3-Manifolds},
preprint (2026).

\bibitem{Caprini2016}
C.~Caprini et al.,
J.\ Cosmol.\ Astropart.\ Phys.\ \textbf{04}, 001 (2016).

\bibitem{ET}
M.~Maggiore et al.,
J.\ Cosmol.\ Astropart.\ Phys.\ \textbf{03}, 050 (2020).

\bibitem{CE}
M.~Evans et al., arXiv:2109.09882.
\end{thebibliography}

\end{document}
